\chapter{시간: 파이프라인과 타이밍 클로저}\label{ch:pipeline}

릴레이식 기계의 시간은 순서기가 정한다. 한 틱이 6.3~ms이고 열 틱이 한 위치다(\ref{sec:t-sequencer}절).
이 6.3~ms는 릴레이가 가장 느릴 때에도 \code{STROBE}, \code{JUDGE},
\code{CLEAR}의 규칙 L1--L9가 지켜지도록 계산해서 고른 값이다(\ref{ch:relay-timing}장). 물음은 하나다. 가장
느린 부품도 정해진 창 안에 일을 마치는가. 마치지 못하면 창을 넓히거나 회로를 바꾼다.

FPGA에도 같은 물음이 있다. \ref{ch:sync}장의 규칙을 다시 적으면, 레지스터 사이의 조합 논리는 한 클럭 주기
안에 끝나야 한다. 다른 점은 창이 밀리초가 아니라 나노초라는 것, 그리고 확인하는 것이 사람의 계산과
시뮬레이션이 아니라 설계 도구라는 것이다. 이 창을 가장 빠듯하게 쓰는 회로는 실물 기계에서 전류가 공짜로
하던 일, 곧 이어진 선을 찾는 회로다(\ref{sec:f-cost}절). 앞 장이 회로의 공간 비용을 다뤘다면 이 장은 시간
비용을 다룬다. 긴 일을 창에 맞게 나누는 방법(파이프라인), 회로가 칩 위에 놓이는 과정(합성, 배치, 배선),
규칙을 지켰는지 확인하는 과정(타이밍 클로저)의 차례로 본다.

\section{파이프라인}\label{sec:pipeline}

\Relay 테스트 레지스터의 처음 회로를 떠올리자(\ref{sec:t-register}절). 시험 전압이 걸려 있는 동안 감지
릴레이 $K$가 붙고, 26개가 모두 붙으면 통과 릴레이 $P$가 붙고, 판정 순간에 $P$가 놓여 있으면 정지 릴레이
$S$가 붙는다. 릴레이 세 단이 한 줄로 이어진 사슬이고, 최악의 시간을 더하면 $5 + 5 + 5 = 15$~ms다. 그런데
기어식에서 드럼 접점이 패드 위에 확실히 얹혀 있는 창은 6~ms였다. 해법은 사슬의 중간에 기억을 끼우는
것이었다. 창 안에서는 $K$가 붙어 자기 유지로 결과를 붙들기만 하고, 판정은 창이 닫힌 뒤에 한다. 그러면 창
안에 들어가야 하는 것은 릴레이 한 단의 동작(최악 5~ms)뿐이다. 긴 사슬을 기억으로 끊어 단마다 짧게 만든
것이다. 대가는 판정이 늦어진다는 것이다. 한 위치가 \code{STROBE}, \code{JUDGE}, \code{CLEAR}의 세 단계를
차례로 거친다(그림~\ref{fig:t-cycle}).

그러나 여러 위치가 이 단계들을 \emph{겹쳐서} 지나지는 않는다. 감지 릴레이는 한 벌뿐이므로, 다음 위치의
\code{STROBE}는 이번 위치의 \code{CLEAR}가 끝난 뒤에야 올 수 있다(L6, 표~\ref{tab:t-rules}). 순서기의 링에서
겹치는 것은 하나뿐이다. 계수기를 다음 위치로 미는 \code{CK}(Q7--Q8)가 \code{CLEAR}(Q6--Q8) 안에 들어 있다
(그림~\ref{fig:t-ring}). 시험 전압이 없는 동안이므로 L9가 이것을 허락한다. 다음 위치의 감지를 이번 위치의
판정과 겹치려면 감지 릴레이가 한 벌 더 있어야 한다. 이 기계는 그렇게 하지 않는다.

\Why FPGA에서 긴 사슬은 조합 논리다. 조합 논리가 길면 신호가 지나는 데 오래 걸리고, 그만큼 클럭을 느리게
해야 한다. 회로에서 가장 긴 경로 하나가 회로 전체의 속도를 정한다. 테스트 레지스터가 감지 릴레이로 사슬을
끊었듯, 긴 조합 논리의 중간에 레지스터를 끼워 넣으면 경로가 둘로 나뉜다. 각각은 원래의 절반쯤만 걸리므로
클럭을 그만큼 빠르게 할 수 있다. 이것이 \term{파이프라인}{pipeline}이다.

대가도 같다. 끼워 넣은 레지스터만큼 플립플롭이 들고, 결과가 나오기까지 클럭이 더 든다(이 지연을
\term{레이턴시}{latency}라 한다). 다른 점은 겹침이다. 감지 릴레이는 \code{CLEAR}라는 단계를 따로 거쳐야
지워지지만(접점 면에서 10~ms 이상), 플립플롭은 지우는 단계 없이 다음 에지에 새 값으로 덮어쓴다. 그래서 앞
단이 한 입력을 다음 단에 넘기고 나면 곧바로 다음 입력을 받을 수 있다. 매 클럭 새 입력을 넣고 매 클럭 결과를
하나씩 받으므로 \emph{처리량}은 줄지 않는다. 세탁기와 건조기를 따로 두는 것과 같다. 빨래 한 통이 끝나는 데
걸리는 시간은 같거나 조금 길어지지만, 세탁이 끝난 통을 건조기로 넘기자마자 다음 통의 세탁을 시작할 수 있다.

\General 수 네 개를 더하는 회로를 보자(그림~\ref{fig:g-pipe}). 위쪽은 덧셈기 세 개가 한 줄로 이어져 있어
신호가 덧셈기 세 개를 연달아 지나야 한다. 아래쪽은 둘씩 먼저 더하도록 순서를 바꾸고 중간에 레지스터를 넣어
두 단으로 나눴다. 첫 단의 덧셈기 둘은 나란히 놓여 동시에 동작하므로, 어느 단에서든 신호가 지나는 덧셈기는
하나다.

\begin{figure}[ht]
\centering
\begin{tikzpicture}[x=1cm,y=1cm,
   add/.style={draw,circle,fill=cbluebg,draw=cblue,minimum size=7mm,inner sep=0pt,font=\small}]
  % 파이프라인 없음
  \node[anchor=west,font=\small\bfseries] at (-0.5,3.6) {파이프라인 없음: 한 클럭에 덧셈 세 번};
  \node[reg] (i0) at (0,2.4) {입력};
  \node[add] (a0) at (2.2,2.4) {+}; \node[add] (a1) at (4.4,2.4) {+}; \node[add] (a2) at (6.6,2.4) {+};
  \node[reg] (o0) at (8.8,2.4) {합};
  \draw[->] (i0) -- (a0); \draw[->] (a0) -- (a1); \draw[->] (a1) -- (a2); \draw[->] (a2) -- (o0);
  \draw[<->,cred] (0.5,1.65) -- node[below,note,text=cred] {덧셈기 세 개의 지연이 클럭 주기를 정한다} (8.3,1.65);
  % 파이프라인 있음
  \node[anchor=west,font=\small\bfseries] at (-0.5,0.3) {두 단 파이프라인: 한 클럭에 덧셈 한 번};
  \node[reg] (i1) at (0,-1.0) {입력};
  \node[add] (b0) at (2.2,-0.5) {+}; \node[add] (b1) at (2.2,-1.5) {+};
  \node[reg,fill=cgreenbg,draw=cgreen] (m) at (4.4,-1.0) {중간};
  \node[add] (b2) at (6.6,-1.0) {+};
  \node[reg] (o1) at (8.8,-1.0) {합};
  \draw[->] (i1) -- (b0); \draw[->] (i1) -- (b1);
  \draw[->] (b0) -- (m); \draw[->] (b1) -- (m); \draw[->] (m) -- (b2); \draw[->] (b2) -- (o1);
  \draw[<->,cgreen] (0.5,-2.25) -- node[below,note,text=cgreen] {1단} (3.9,-2.25);
  \draw[<->,cgreen] (4.9,-2.25) -- node[below,note,text=cgreen] {2단} (8.3,-2.25);
  \node[note,anchor=west,text width=40mm,align=left] at (10.0,-1.0) {초록 레지스터가 추가된 것. 각 단의 지연이 짧아져 클럭을 빠르게 할 수 있다.};
\end{tikzpicture}

\medskip
\begin{tikztimingtable}[timing/slope=0.1,timing/coldist=4pt,xscale=2.4,yscale=1.25,semithick]
  \code{clk} & 7{.5L .5H} \\
  입력       & .5D{A} 1D{B} 1D{C} 1D{D} 1D{E} 1D{F} 1D{G} .5D{H} \\
  중간       & .5U 1D{A} 1D{B} 1D{C} 1D{D} 1D{E} 1D{F} .5D{G} \\
  합         & 1.5U 1D{A} 1D{B} 1D{C} 1D{D} 1D{E} .5D{F} \\
\extracode
  \begin{pgfonlayer}{background}
    \vertlines[cgray!50,dashed]{0.5,1.5,...,6.5}
  \end{pgfonlayer}
\end{tikztimingtable}
\caption{덧셈 회로의 파이프라인. 아래 파형은 두 단 파이프라인의 동작이다. 글자는 입력 묶음의 이름이다.
결과는 두 클럭 늦게 나오지만 매 클럭 하나씩 나온다.}
\label{fig:g-pipe}
\end{figure}

\Ours FPGA 엔진에서 스크램블러의 안쪽 치환을 계산하는 부분이 파이프라인이다. 글자 하나가 가운데 로터,
왼쪽 로터, (그리스 로터,) 반사판을 지나 다시 돌아오는 계산인데, 이것을 한 클럭에 다 하면 표 조회 일곱
번이 한 줄로 이어진 긴 경로가 된다. 그래서 로터 한 단을 지날 때마다 레지스터를 두어 일곱 단으로 나눴다.
글자 26개를 차례로 흘려 넣으면 26개의 결과가 차례로 나온다. 릴레이 기계와 달리 여기서는 단마다 다른 글자가
동시에 들어 있다.

처음의 일곱 단 구성은 그래도 목표 클럭을 맞추지 못했다. 로터 한 단의 통과는 사실 두 가지 일이다.
글자에 회전량을 더해 표의 주소를 만드는 일($x + s \bmod 26$)과, 표를 읽고 회전량을 빼는 일이다.
두 일을 한 단에 넣으면 덧셈, 26으로 나눈 나머지, 표 조회, 뺄셈이 한 클럭 안에 이어진다. 각 단을 이 두
일로 다시 쪼개 열네 단으로 만들자 타이밍이 맞았다(그림~\ref{fig:g-inner-pipe}).

\begin{figure}[ht]
\centering
\begin{tikzpicture}[x=1cm,y=1cm]
  \node[anchor=west,font=\small\bfseries] at (-0.3,2.2) {7단: 로터 한 단 통과가 한 클럭};
  \foreach \i/\t in {0/M,1/L,2/G,3/반사,4/G,5/L,6/M} {
    \node[blkb,minimum width=13mm,minimum height=8mm,font=\footnotesize] (p\i) at (\i*1.75+0.5,1.3) {\t};
    \draw[fill=cyellowbg] (\i*1.75+1.3,0.9) rectangle ++(0.22,0.8);
  }
  \foreach \i/\j in {0/1,1/2,2/3,3/4,4/5,5/6} { \draw[->] (\i*1.75+1.52,1.3) -- (p\j); }
  \draw[->] (-0.6,1.3) -- (p0);
  \node[anchor=west,font=\small\bfseries] at (-0.3,0.0) {14단: 한 단을 둘로 (주소 만들기 / 표 읽기)};
  \foreach \i in {0,...,13} {
    \pgfmathtruncatemacro{\odd}{mod(\i,2)}
    \ifnum\odd=0
      \node[blkb,minimum width=5.5mm,minimum height=8mm,font=\tiny,inner sep=1pt] (q\i) at (\i*0.875+0.1,-0.9) {주소};
    \else
      \node[blkg,minimum width=5.5mm,minimum height=8mm,font=\tiny,inner sep=1pt] (q\i) at (\i*0.875+0.1,-0.9) {읽기};
    \fi
    \draw[fill=cyellowbg] (\i*0.875+0.48,-1.3) rectangle ++(0.14,0.8);
  }
  \node[note,anchor=west,text width=30mm,align=left] at (12.2,1.3) {노란 막대가\\레지스터};
  \node[note,anchor=west,text width=30mm,align=left] at (12.2,-0.9) {단마다 할 일이\\절반};
\end{tikzpicture}
\caption{안쪽 치환 파이프라인의 분할. M, L, G는 가운데, 왼쪽, 그리스 로터의 통과다. 단 수가 두 배가 되어
결과가 나오기까지 클럭이 더 들지만, 이 계산은 가운데 로터가 바뀔 때에만 하므로 전체 시간에는 거의 영향이
없다.}
\label{fig:g-inner-pipe}
\end{figure}

이 예에서 볼 점은 파이프라인을 \emph{어디에} 썼는가다. 안쪽 치환의 계산은 26위치에 한 번만 하는 일이므로
레이턴시가 늘어도 아깝지 않다. 반면 전파 논리(상태 레지스터에서 나와 다시 상태 레지스터로 들어가는 고리)는
단순히 레지스터를 끼워 나눌 수 없다. 한 걸음의 결과가 다음 걸음의 입력이므로, 한 걸음을 두 클럭으로 나누면
걸음마다 두 클럭이 든다. 릴레이 기계에는 이런 걸음이 없다. 전류가 한 \code{STROBE} 안에 덩어리 전체로
퍼지기 때문이다(\ref{sec:f-current}절). FPGA는 그 퍼짐을 걸음으로 나눠 클럭마다 하나씩 하고
(\ref{sec:f-rounds}절), 그 한 걸음은 다시 한 클럭 안에 들어가야 한다. 그래서 그쪽은 논리 자체를 짧게 만드는
방법으로 풀어야 한다(\ref{ch:lesson}장).

표~\ref{tab:p-stages}는 두 기계의 단계 나누기를 나란히 놓은 것이다.

\begin{table}[ht]
\centering
\caption{테스트 레지스터의 세 단계와 FPGA의 파이프라인.}
\label{tab:p-stages}
\small
\begin{tabularx}{\linewidth}{>{\raggedright\arraybackslash}p{30mm}>{\raggedright\arraybackslash}X>{\raggedright\arraybackslash}X}
\toprule
& 릴레이 테스트 레지스터 & FPGA 파이프라인 \\
\midrule
사슬을 끊는 기억 & 감지 릴레이 26개(자기 유지) & 레지스터(플립플롭) \\
단을 나누는 것 & 순서기의 창 \code{STROBE}, \code{JUDGE}, \code{CLEAR} & 클럭 에지 \\
한 단에 들어갈 것 & 릴레이 한 단의 동작(최악 5~ms) & 두 레지스터 사이의 조합 논리(한 주기) \\
기억을 지우는 일 & \code{CLEAR}라는 단계가 따로 있다 & 다음 에지에 덮어쓴다 \\
다음 일과 겹치는가 & 겹치지 않는다(\code{CK}만 \code{CLEAR} 안에) & 매 클럭 새 입력이 들어간다 \\
\bottomrule
\end{tabularx}
\end{table}

\section{합성, 배치, 배선}

\Relay 릴레이 보드도 회로에서 기판까지 여러 단계를 거친다. 회로도 생성기가 ``무엇이 무엇에 이어지는가''를
정하고, 기판 생성기가 그 넷리스트에서 출발해 부품의 자리를 공식으로 정하고(\ref{sec:gen-place}절), 동선은
사람이 KiCad에서 손으로 긋는다(\ref{sec:gen-noroute}절). 이 동선의 길이는 기계의 시간에 거의 영향이 없다.
전류가 동선을 지나는 시간은 릴레이가 넘어가는 수 밀리초에 비하면 없는 것과 같다(\ref{sec:f-current}절).
배선이 정하는 것은 오히려 크기다. 시프터 보드는 16칸 단에서 선이 위아래로 최대 16칸을 건너가 수직 배선이
20개 넘게 겹친다. 단 사이에 배선 통로를 남기느라 보드가 처음 가정한 160 $\times$ 200~mm에 들어가지 않고
250 $\times$ 228~mm가 되었다(그림~\ref{fig:g-shifter-pcb}).

FPGA에서도 지금까지 ``합성 도구가 코드를 회로로 바꾼다''고만 했다. 실제로는 세 단계이고, 세 단계를 모두
도구가 한다(그림~\ref{fig:g-flow}).

\begin{figure}[ht]
\centering
\begin{tikzpicture}[x=1cm,y=1cm]
  \node[blkn,minimum width=22mm,minimum height=12mm] (hdl) at (0,0) {HDL 코드};
  \node[blkb,minimum width=24mm,minimum height=12mm] (syn) at (3.3,0) {합성\\{\footnotesize synthesis}};
  \node[blkb,minimum width=24mm,minimum height=12mm] (pl) at (6.8,0) {배치\\{\footnotesize place}};
  \node[blkb,minimum width=24mm,minimum height=12mm] (rt) at (10.3,0) {배선\\{\footnotesize route}};
  \node[blkg,minimum width=22mm,minimum height=12mm] (bit) at (13.6,0) {비트스트림};
  \draw[->] (hdl) -- (syn); \draw[->] (syn) -- (pl); \draw[->] (pl) -- (rt); \draw[->] (rt) -- (bit);
  \node[note,below=2mm of syn] {논리를 LUT와\\플립플롭으로};
  \node[note,below=2mm of pl] {각 LUT를 칩의\\어느 자리에};
  \node[note,below=2mm of rt] {자리 사이를\\어느 배선으로};
  \node[note,below=2mm of bit] {칩에 써넣는\\설정 파일};
  \draw[decorate,decoration={brace,amplitude=5pt}] (2.0,0.9) -- (4.6,0.9) node[midway,above=6pt,note] {크기가 정해진다};
  \draw[decorate,decoration={brace,amplitude=5pt}] (5.5,0.9) -- (11.6,0.9) node[midway,above=6pt,note] {지연이 정해진다};
\end{tikzpicture}
\caption{HDL 코드가 FPGA에 올라가기까지.}
\label{fig:g-flow}
\end{figure}

\Why 릴레이 보드와 달리 FPGA에서는 배선이 시간을 정한다. FPGA에서 신호의 지연은 두 부분으로 이루어진다.
LUT를 지나는 시간과 LUT 사이의 \emph{배선}을 지나는 시간이다. 배선의 지연은 두 LUT가 칩 위에서 얼마나
떨어져 놓였는지, 그 사이를 어떤 배선으로 이었는지에 달려 있고, LUT의 지연보다 큰 경우가 흔하다. 그러므로
회로의 크기는 합성만 하면 알 수 있지만, 회로의 속도는 배치와 배선까지 끝나야 정확히 알 수 있다. 합성 직후에
도구가 알려 주는 속도는 배선의 지연을 어림한 \emph{추정값}이다. 그리고 놓을 것이 LUT 수만 개이므로, 자리와
길을 사람이 고르지 않고 도구가 자동으로 고른다. 릴레이 보드에서 사람이 하는 배선을 FPGA에서는 도구가 하고,
그 결과가 곧 속도가 된다.

\General 플립플롭 두 개 사이에 LUT 하나가 있는 회로도, 도구가 두 플립플롭을 칩의 양 끝에 놓으면 느려진다.
회로가 작을 때에는 도구가 관련된 부품을 가까이 모아 놓을 수 있지만, 칩이 꽉 찰수록 놓을 자리가 없어 부품이
흩어지고 배선이 길어진다. 그래서 같은 설계도 칩을 절반쯤 채웠을 때보다 거의 다 채웠을 때 속도를 맞추기가
훨씬 어렵다. 시프터 보드에서 배선 통로가 모자라 보드가 커진 것과 같은 일이다. 다만 칩의 크기는 정해져
있으므로, FPGA에서는 커지는 대신 느려진다.

\Ours 이 프로젝트의 도구는 Xilinx Vivado다. FPGA 엔진의 크기(LUT 41,040개, 대상 칩의 30.5\%)는 합성
단계의 결과다. 엔진을 여러 대 넣으면 칩이 차오르므로, 한 대일 때 맞던 타이밍이 여러 대일 때에도 맞는지는
배치와 배선을 해 보아야 안다. 보드별 엔진 수는 그 결과를 보고 정한다.

\section{타이밍 클로저}\label{sec:timing}

\Relay 릴레이 기계에서 시간을 확인하는 일은 세 걸음이었다. 먼저 테스트 레지스터 안의 접점이 지켜야 할
조건을 릴레이의 최악 시간으로 적는다. 예를 들어 \code{STROBE}의 접점은 5~ms 이상 닫혀야 한다(접점 면).
다음으로 커넥터 핀과 그 접점 사이에 있는 타이밍 릴레이의 지연 범위를 가장 나쁘게 넣어, 뱅크가 지킬 규칙
L1--L9로 옮긴다(커넥터 면, \ref{sec:t-faces}절). 마지막으로 순서기의 박자를 그 규칙에 맞춘다. 한 틱 6.3~ms는
L1의 7.5~ms에 순서기 자신의 출력 릴레이가 깎아 먹는 폭과 여유를 더해 정한 값이다. 시뮬레이션이 릴레이 시간을
가장 나쁘게 넣어 잰 \code{STROBE}의 최소 폭은 8.23~ms로, 규칙보다 0.73~ms 길었다(\ref{sec:t-sequencer}절). 이
0.73~ms가 이 절에서 \emph{슬랙}이라 부를 값이다. 확인은 시뮬레이터가 한다. 릴레이 시간을 가장 느리게, 가장
빠르게, 또는 섞어서 넣고, 여러 운전 시나리오와 17,576위치 전체를 돌려 커넥터 핀에서 L1--L9를 잰다
(\ref{sec:v-time}절).

\Why FPGA에서는 배선까지 끝나면 도구가 회로의 모든 레지스터 사이 경로에 대해 지연을 계산하고, 클럭 주기와
비교한다. 주기에서 지연을 뺀 값을 \term{슬랙}{slack, 여유}이라 한다. 슬랙이 양수이면 신호가 제때 도착하고,
음수이면 늦는다. 가장 나쁜 경로의 슬랙이 0 이상이 되도록 설계를 고쳐 나가는 과정이
\term{타이밍 클로저}{timing closure}다. 경로 하나라도 슬랙이 음수이면 그 회로는 그 클럭에서 동작을 보장할
수 없다. ``대체로 동작한다''는 없다. 온도나 전압에 따라 가끔 틀린 값을 내는 회로가 되기 때문이다.

생각의 틀은 릴레이와 같다. 가장 나쁜 지연이 정해진 창에 들어가는지 보고, 남는 시간을 슬랙으로 적는다.
느리게 하면 안전하다는 것도 같다. 클럭 주기를 늘리면 모든 경로의 슬랙이 그만큼 커진다. 다른 점은 셋이다.

\begin{itemize}
\item \textbf{무엇이 정해져 있는가.} 릴레이 기계에서는 부품의 지연이 넓게 흩어져 있고 최소를 모르는 것이
많다. 그래서 규칙을 하한으로 적고 박자 쪽을 넉넉하게 정했다(\ref{sec:t-bounds}절). 박자는 운전 중에도
바뀐다. 관찰 모드로 바꾸든 한 걸음씩 누르든, 느려지는 쪽이면 언제든 괜찮다. FPGA에서는 클럭 주기가 설계
전에 정한 숫자 하나이고 운전 중에 바뀌지 않는다. 바뀌는 것은 회로 쪽이다. 코드를 고치거나 배치가 달라질
때마다 경로의 지연이 달라지고, 그때마다 모든 경로가 그 주기 안에 드는지 다시 확인한다.
\item \textbf{어떻게 확인하는가.} 릴레이의 시뮬레이터는 회로를 \emph{돌려 보고} 핀에서 잰다. 돌려 본
시나리오만 확인된다. FPGA 도구는 회로를 돌리지 않는다. 칩의 지연 모형으로 경로마다 지연을 \emph{계산}하므로,
어떤 입력이 오든 모든 경로가 확인된다. 이것을 \term{정적 타이밍 분석}{static timing analysis}이라 한다.
\item \textbf{너무 빠른 쪽.} FPGA 도구는 반대쪽도 본다. 한 에지에서 앞 레지스터가 내보낸 새 값이 너무 빨리
도착하면, 다음 레지스터가 같은 에지에서 옛 값을 받아 적기 전에 그 값을 덮어 버린다. 값이 한 클럭에 레지스터
두 개를 건너는 셈이다. 이 검사를 \term{홀드}{hold} 검사라 하고, 앞에서 말한 ``제때 도착하는가''의 검사를
셋업 검사라 한다(\ref{sec:t-faces}절의 셋업 시간과 같은 말이다). 홀드 위반은 클럭을 늦춰도 고쳐지지 않는다.
같은 에지 안에서 벌어지는 경주이기 때문이다. 릴레이 계수기는 같은 위험을 구조로 막았다. 마스터와 슬레이브의
문이 동시에 열리지 않게 해서 값이 한 클럭에 두 번 넘어가지 않게 했다(\ref{sec:relay-counter}절). FPGA에서는
경로마다 지연이 다르므로 도구가 이것도 숫자로 잰다.
\end{itemize}

\General 클럭이 100MHz이면 주기는 10나노초다(그림~\ref{fig:g-slack}). 어떤 경로의 지연이 8나노초이면
슬랙은 $+2$나노초로 여유가 있다. 12나노초이면 슬랙은 $-2$나노초로 실패다. 실패한 경로를 고치는 방법은
세 가지다. 그 경로를 파이프라인으로 나누거나, 논리의 구조를 바꿔 거치는 LUT의 수(논리 단수)를 줄이거나,
클럭을 느리게 한다. 테스트 레지스터의 처음 회로가 창에 들어가지 않았을 때에도 같은 갈림길이 있었다. 기어를
늦추면 들어가지만 최악값을 담으려면 속도가 3분의 1이 되었다. 그래서 늦추는 대신 사슬을 감지 릴레이로 끊는
쪽, 곧 파이프라인에 해당하는 쪽을 골랐다(\ref{sec:t-register}절).

\begin{figure}[ht]
\centering
\begin{tikzpicture}[x=0.9cm,y=1cm]
  \draw[->] (0,0) -- (14,0) node[right,font=\small] {시간 (ns)};
  \foreach \t in {0,2,4,6,8,10,12} { \draw (\t,0.08) -- (\t,-0.08) node[below,font=\footnotesize] {\t}; }
  \draw[cblue,line width=1.2pt,dashed] (0,-0.5) -- (0,3.3) node[above,font=\footnotesize] {클럭 에지};
  \draw[cblue,line width=1.2pt,dashed] (10,-0.5) -- (10,3.3) node[above,font=\footnotesize] {다음 에지};
  \fill[cgreen!60] (0,2.2) rectangle (8,2.8);
  \node[font=\footnotesize,anchor=west] at (0.2,2.5) {경로 A: 지연 8 ns};
  \draw[<->,cgreen] (8,2.5) -- node[above,font=\footnotesize,text=cgreen] {슬랙 $+2$} (10,2.5);
  \fill[cred!55] (0,1.0) rectangle (12,1.6);
  \node[font=\footnotesize,anchor=west] at (0.2,1.3) {경로 B: 지연 12 ns};
  \draw[<->,cred] (10,1.9) -- (12,1.9) node[right,font=\footnotesize,text=cred] {슬랙 $-2$ (실패)};
\end{tikzpicture}
\caption{슬랙. 100MHz 클럭(주기 10ns)에서의 두 경로.}
\label{fig:g-slack}
\end{figure}

\Ours FPGA 엔진의 목표 클럭은 처음에 보드의 클럭과 같은 100MHz였다. 전파 회로를 케이블 단위 구조로
바꾼 뒤(\ref{ch:lesson}장) 합성 단계에서
추정한 슬랙은 $+0.18$나노초였고, 가장 느린 경로는 전파 논리로 LUT를 11단 거친다. 여유가 0.18나노초라는
것은 거의 없다는 뜻이고, 더구나 배선 지연을 어림한 추정값이다. 실제로 배치를 마치자 슬랙이 $-3.4$나노초,
배선 도중에는 $-5.2$나노초가 되었다. 전파 경로가 약 15나노초 걸렸고, 그 대부분이 LUT가 아니라 배선의
지연이었다. 676비트의 상태가 메뉴의 선 16개의 양쪽 끝, 곧 32곳의 선택 회로로 퍼졌다가 다시 한 레지스터로
모여야 하므로 칩 위에서 먼 거리를 가는 선이 많다.

가장 느린 경로가 전파 회로라는 것은 이 문서의 줄거리와 맞는다. 실물 기계에서 전류가 공짜로 하던 일, 곧
이어진 선을 한 걸음 퍼뜨리는 일이 FPGA에서는 676비트 전체를 읽어 676비트 전체를 다시 쓰는 회로가
되었다(\ref{sec:f-rounds}절). 그 회로가 칩 위에 넓게 퍼져야 했고, 그 거리가 클럭을 정했다.

고칠 방법은 앞의 셋이다. 테스트 레지스터는 늦추는 대신 구조를 바꿨지만, FPGA 엔진은 반대로 클럭을 늦췄다.
전파 회로는 고리라서 파이프라인으로 나누면 얻는 것이 없고, 구조를 더 바꾸는 것은 큰 일이다. 이 프로젝트의
최종 목표는 실물이고 FPGA는 세 번째 전시물이자 기준 장치이므로, 시스템 클럭을 50MHz로 내렸다. 50MHz에서는
엔진 한 대일 때 슬랙이 $+2.40$나노초, 두 대일 때 $+0.13$나노초다. 회로는 똑같은 두 벌인데 칩이 차서 배선이
길어진 것이다. 자세한 경과는 \ref{sec:e-results}절에 있다. 늦춘 뒤에도 FPGA는 위치 하나에 약 0.3마이크로초로,
실물 두 기계보다 약 10만 배 빠르다(표~\ref{tab:f-speed}).

\begin{keybox}[title={회로 설계의 두 가지 실패}]
회로 설계가 실패하는 방식은 두 가지다. 회로가 칩에 \emph{들어가지 않거나}(자원), 원하는 클럭에서
\emph{제때 끝나지 않거나}(타이밍). 프로그램은 느려도 답은 맞게 내지만, 타이밍을 맞추지 못한 회로는 답이
틀린다. 릴레이 기계도 같다. \code{JUDGE}가 짧아 정지 릴레이가 붙지 못하면(L3) 정지를 놓친다. 그래서 회로
설계에서는 ``동작한다''와 ``들어간다''와 ``제 속도로 돈다''를 따로 확인한다.
\end{keybox}

\section{정리}

\begin{itemize}
\item 릴레이 기계와 FPGA는 같은 물음을 다룬다. 가장 느린 부품이나 경로도 정해진 창 안에 일을 마치는가. 창은
릴레이에서는 순서기의 틱(6.3~ms), FPGA에서는 클럭 주기(나노초)다.
\item 긴 사슬은 중간에 기억을 끼워 나눈다. 테스트 레지스터는 감지 릴레이로, FPGA는 레지스터로 나눈다.
릴레이 기계는 감지 릴레이가 한 벌이고 지우는 단계가 따로 있어 위치들을 겹치지 않는다. FPGA의 파이프라인은
매 클럭 새 입력을 받는다.
\item 릴레이 보드에서 배선은 크기를 정하고 시간에는 거의 영향이 없다. FPGA에서는 도구가 배치와 배선을 하고,
배선의 지연이 속도를 정한다.
\item 도구는 모든 경로의 슬랙을 계산한다. 늦게 도착하는 쪽(셋업)은 클럭을 늦추면 풀리지만, 너무 일찍
도착하는 쪽(홀드)은 풀리지 않는다.
\item 이 프로젝트에서 가장 느린 경로는 전파 회로였다. 전류가 공짜로 하던 일이 FPGA에서는 클럭을 정하는
회로가 되었고, 엔진은 100MHz 대신 50MHz로 돈다.
\end{itemize}
